\documentclass[11pt]{article}
\input{../MilestoneDraft.text}
\input{../Common.text}
\title{Development Plan\\\progname}
\author{\authname}
\date{September 23, 2026}
\begin{document}
\maketitle
\draftnotice
Table~\ref{tab:dev-revisions} summarizes revision history.
\begin{center}
\begin{minipage}{\textwidth}
\captionof{table}{Revision history.}\label{tab:dev-revisions}
\begin{tabularx}{\textwidth}{llX}
\toprule Date & Contributor & Change\\\midrule
2026-09-23 & Aaron Loh / AI assistance & Initial review draft based on Aaron's planning interview; team ratification pending.\\\bottomrule
\end{tabularx}
\end{minipage}
\end{center}

This plan describes how a four-person team proposes to build and evaluate a
general-purpose, voice-supervised agent interface for blind users. It covers
team operations, technical choices, the proof of concept (POC), and evaluation.
It does not promise a universal replacement for screen readers. Initial use
is on Mac, in English, alongside VoiceOver and keyboard access. The companion
Problem Statement and Goals defines the user outcomes.

\section{Confidential Information?}
No external confidentiality agreement is known from Aaron's planning interview.
\pending{All members must confirm any employer, sponsor, or industry restrictions.}
Participant identities, recordings, account credentials, and personal task data
are private even without an industry agreement. Public artifacts will use test
accounts, synthetic data, and consented, de-identified findings. Raw recruitment
and meeting materials remain outside the public repository. No raw microphone
audio will be retained by default; provider retention policies require separate
review and disclosure.

\section{Intellectual Property (IP) to Protect}
No external intellectual-property agreement is currently known. The intended
contribution is the accessible interaction and supervision layer, its adapters,
and evaluation, building on existing models and tools. The team must confirm
ownership and any required institutional agreements before release. Upstream
code, model weights, services, and assets retain their own terms; licensing our
code does not relicense dependencies.

\section{Copyright License}
MIT is the proposed license for original project contributions. Its permissive
terms require retaining the copyright and permission notice \citep{MITLicense}.
The repository's \href{https://github.com/NaturalControl/natural-control/blob/main/LICENSE}{LICENSE}
currently contains the inherited template notice. The local license proposal
preserves that notice; team ratification and confirmation of attribution remain
pending before publication. Third-party notices will be recorded when actual
dependencies are selected, rather than assuming all candidate tools share MIT.

\section{Team Meeting Plan}
The proposal is one weekly team sync, supported by asynchronous work.
\pending{Recurring day, time, duration, and virtual/in-person arrangement.}
Each meeting has a rotating chair and note taker, an issue-based agenda, attendance,
decisions, owners, and next commitments recorded in GitHub. Members review
materials beforehand. A member unable to attend can provide an advance written
update and input on decisions; emergencies are communicated when practical.
A supervisor/teaching assistant (TA) check-in approximately monthly is proposed, in addition to
required course reviews; the TA's availability and actual schedule need confirmation.

\section{Team Communication Plan}
Discord is proposed for coordination, with a response to actionable messages
within one working day. Decisions, acceptance criteria, blockers, and assignments
move to GitHub issues so they remain discoverable. Sensitive participant details
and secrets do not belong in either public issues or the project board. Urgent
blockers are explicitly flagged; silence is not agreement or approval.

\section{Team Member Roles}
The team has four members; only Aaron Loh's name is confirmed in this draft.
Primary ownership provides accountability without preventing another member from
contributing. All four members implement, test, review, and document work.
Owners and cross-review assignments are revisited at milestones.

Table~\ref{tab:dev-roles} summarizes proposed workstreams.
\begin{center}
\begin{minipage}{\textwidth}
\captionof{table}{Proposed workstreams.}\label{tab:dev-roles}
\begin{tabularx}{\textwidth}{>{\raggedright\arraybackslash}p{0.29\textwidth}X}
\toprule Proposed workstream & Responsibility and assignment\\\midrule
Voice and accessible client & Conversation controls, keyboard access, VoiceOver compatibility. Owner and cross-reviewer pending.\\
Agent and tool integration & Hermes adapter, account/browser tools, local execution control. Owner and cross-reviewer pending.\\
Deployment and reliability & Configuration, credentials, CI, reconnection, observability. Owner and cross-reviewer pending.\\
Evaluation and quality & Recruitment coordination, task protocol, outcome measurement, accessibility checks. Owner and cross-reviewer pending.\\\bottomrule
\end{tabularx}
\end{minipage}
\end{center}
Meeting chair, note taker, release coordinator, and TA liaison are rotating process
roles. A contributor reviews outside their primary workstream to reduce dependence
on one person. \pending{Three names, named initial owners, cross-reviewers, and liaison.}

\section{Workflow Plan}
Work is tracked in \href{https://github.com/NaturalControl/natural-control}{the public repository}.
Each meaningful change starts with an issue describing the problem, acceptance
criteria, owner, and relevant risk. Labels distinguish features, defects,
accessibility, research, documentation, and blockers. Short-lived topic branches
and small pull requests keep changes reviewable. One non-author reviewer must
approve each pull request; relevant continuous integration (CI) checks must pass before merging.

The repository already contains a LaTeX/PDF and Pages workflow. Python and
TypeScript checks will be added when the corresponding implementation exists:
format/lint, type checks, unit/contract tests, and applicable integration tests.
Passing automation does not establish accessibility. Manual VoiceOver/keyboard
checks and participant evidence remain necessary. Secrets stay out of source
control and logs. A change is done when its acceptance criteria, checks,
documentation, and review are complete, with known limitations recorded.

\section{Project Decomposition and Scheduling}
A GitHub Project is proposed with Backlog, Ready, In progress, Review, and Done
columns. Each member commits to concrete weekly outputs rather than fixed hours
or commit quotas. Dependencies and risks are recorded on issues. The team aims
for one primary active item per member and resolves blocked work at the weekly
sync. \pending{Create and link the actual project board; agree initial assignments.}

Drafts are ready 48 hours before deadlines. For the first milestone, the target
is September 26 at 11:59 p.m. Toronto time, leaving time for human review and
ratification before September 28 at 11:59 p.m. The TA's September 23 guidance was
that assessment uses the repository's latest push at the deadline, without an
Avenue upload; the knowledge base links the meeting source. Later official
announcements take precedence.

Table~\ref{tab:dev-schedule} summarizes milestone overview.
\begin{center}
\begin{minipage}{\textwidth}
\captionof{table}{Milestone overview.}\label{tab:dev-schedule}
\begin{tabularx}{\textwidth}{>{\raggedright\arraybackslash}p{0.25\textwidth}X}
\toprule Period/date & Planned output or course milestone\\\midrule
September 28, 2026 & Problem Statement and Goals; Development Plan including POC and charter.\\
October 1 & First milestone peer-review issues; substantive findings, not manufactured issue counts.\\
October 19 & Software Requirements Specification (SRS) and Hazard Analysis; discovery informs scope and safety requirements.\\
October--November & Working voice/agent supervision loop; instrumented POC trials; repeated discovery if recruitment succeeds.\\
November 2 & Verification and Validation (V\&V) Plan.\\
November 16 & Design revision minus one and proposed fall Interview Report, subject to instructor approval.\\
November 23--December 3 & Course POC demonstration window; assigned slot pending.\\
January 18, 2027 & Design documentation; integration and accessibility refinement continue.\\
March 8 & V\&V report and proposed winter Usability Report, subject to approval.\\
March 31 & Final documentation.\\
April 6 & EXPO; event details to confirm.\\\bottomrule
\end{tabularx}
\end{minipage}
\end{center}
Dates reflect the checked course calendar snapshot \citep{CourseCalendar}; this
is a milestone overview, not a replacement for the complete course calendar.

\paragraph{Resources and budget.}
The proposed total budget has a Canadian-dollar ceiling of \$500, not authorization to spend. A provisional
allocation is \$240 participant appreciation, \$120 model/voice application programming interfaces (APIs), \$60
hosting/telephony, and \$80 reserve. Appreciation assumes at most twelve paid
sessions at \$20, subject to recruitment and institutional guidance. Team Macs
are already available. Test accounts and synthetic data reduce exposure of
personal information. Usage metering and monthly reconciliation will track
spend; paid expansion requires a revised allocation within the ceiling or an
explicitly approved budget change.

\section{Proof of Concept Demonstration Plan}
The POC tests a narrow integration risk before broader prototype development.
\subsection{Risk and experiment}
The central risk is whether a blind user can understand and control a running
general-purpose agent through voice without visually monitoring it. Individual
speech and agent components do not establish that their integration works.
The POC will test the supervision interface across three task classes, rather
than implement a set of website-specific scripts:
\begin{enumerate}
\item \textbf{Email/calendar tools:} use synthetic correspondence to prepare a
reply or calendar entry; clarify ambiguous intent and obtain approval before
sending or creating the event.
\item \textbf{An unfamiliar website:} locate and compare information through
general browser tools, report a meaningful choice, and incorporate a correction.
The site and task are selected before the trial, without coding a dedicated workflow.
\item \textbf{A local Mac task:} locate and open a test file or interact with a
supported application; prepare a derived file, confirming consequential changes
such as overwrite. Native control support is an experimental dependency.
\end{enumerate}

Each class exercises clarification, consequential-action approval, progress,
correction, explicit execution stop, and reconnection. The demonstration uses
controlled accounts and data; no participant's real email is required.

\subsection{Control semantics}
Ordinary speech interruption stops spoken output while the underlying task
continues. An explicit ``stop working'' requests cancellation of execution.
The interface first acknowledges receipt, then separately confirms whether
execution has stopped. Acknowledgment is not proof of cancellation. Once stopped
is confirmed, no further action may be dispatched for that task. An action
already completed cannot be undone by saying stop; in-flight actions with an
unknown outcome must be reported and reconciled before retrying.

Disconnection pauses new actions unless the user previously authorized background
continuation for that task. Reconnection presents current state and pending
decisions; it must not duplicate a send or other consequential action. Navigation
and preparation can proceed autonomously, but sending, purchasing, submitting,
and consequential account/file changes require specific approval. If the action
or its relevant details change, the previous approval does not apply.

The POC will inject a misunderstanding, interrupt speech, request stop during a
tool operation, and disconnect while work is active. Logs correlate task, action,
approval, cancellation, and connection events without retaining sensitive payloads.
The team must verify cancellation through the adapter and downstream worker;
a browser or native action may have different cancellation behavior.

\subsection{Measures and participant evaluation}
Targets below are provisional goals, not results:
\begin{itemize}
\item At least 80\% correct independent completion on a declared representative
task set, following assisted setup. Correct means the specified outcome was
achieved and verified without researcher task assistance. Report numerator,
denominator, task categories, failures, and any assistance separately.
\item Stop acknowledgment within two seconds, measured from the end of a
recognized explicit stop utterance to audible acknowledgment. Record recognition
failures separately and measure execution-stop confirmation latency independently;
no two-second execution-cancellation guarantee is claimed.
\item Approval before every tested consequential action, checked against action
and approval traces, including attempts that should be blocked.
\item Assess whether users can identify completion or a blocker and correct a
seeded misunderstanding. Record errors, clarification turns, time, perceived
control, and useful versus excessive spoken feedback.
\end{itemize}
The team seeks four to six blind participants with screen-reader experience for
repeated discovery and evaluation. Recruitment is unconfirmed. The protocol,
consent and institutional research requirements must be resolved before sessions.
A personal route to an accessibility organization is a possible introduction,
not a partnership or recruitment commitment.

Before the main evaluation, freeze task definitions, outcome checks, allowed
assistance, and exclusions. Report all attempted tasks and explain invalidated
trials. Compare participants' usual screen-reader workflow with the assistant
where feasible, varying order to reduce learning effects. A small convenience
sample supports formative findings, not population-wide efficacy claims.
Technical POC results and participant results will be reported separately.

\subsection{Failure consequences}
Table~\ref{tab:dev-risks} summarizes pOC uncertainties and failure consequences.
\begin{center}
\begin{minipage}{\textwidth}
\captionof{table}{POC uncertainties and failure consequences.}\label{tab:dev-risks}
\begin{tabularx}{\textwidth}{>{\raggedright\arraybackslash}p{0.25\textwidth}X}
\toprule Uncertainty & Evidence needed and consequence if unresolved\\\midrule
Harness supervision & Verify streamed status, approvals, steering, cancellation, and reconnect on a pinned Hermes version. If unavailable, revisit the adapter/harness before promising supervision.\\
Tool cancellation & Test slow and in-flight account/browser operations. If outcome is unknown, expose that state and block retries; do not claim confirmed stop.\\
Native Mac control & Prove permissions, actionable targets, and result verification. If unsuccessful, defer native control and retain general browser/account tasks.\\
Voice and VoiceOver & Measure latency, intelligibility, and competing audio with keyboard alternatives. If remote calls fail, retain local voice.\\
Remote continuity & Test network loss, pending approvals, and duplicate prevention. Defer telephone interaction if it cannot preserve control.\\
Recruitment & Confirm access early. Without participants, report technical feasibility only and disclose the missing usability evidence.\\\bottomrule
\end{tabularx}
\end{minipage}
\end{center}

\section{Expected Technology}
The proposed architecture is:
\begin{quote}
Accessible TypeScript/React web client $\rightarrow$ LiveKit/Python integration
$\rightarrow$ Hermes $\rightarrow$ account tools, Browser Use, and local Mac tools.
\end{quote}
The team will implement the accessible task/decision presentation and supervision
adapter, reusing existing speech, agent, and tool implementations. One supported
default configuration will be documented and version-pinned, with extension
points and user-provided credentials. Provider choices and exact versions remain
pending POC evidence. No model training is planned.

\subsection{Agent and action tools}
\textbf{Hermes is the primary harness candidate.} Its documented programmatic
interface includes streaming events, steering, interrupt requests, and server
requests for clarification/approval \citep{HermesIntegration}. These are promising
integration primitives, not proof that our end-to-end control semantics already
work. The POC must verify actual behavior and failure handling on a pinned release.
Codex app server is an alternative with a custom-client integration surface
\citep{CodexServer}; OpenClicky is an implementation reference. Neither establishes
that arbitrary desktop plugins or account integrations transfer to our client.
OpenClaw is excluded from the proposed stack.

\textbf{Browser Use} is the intended general browser component through Hermes.
Hermes documents a Browser Use path and fallback browser mechanisms
\citep{HermesBrowser}. Packaged browser sessions must not be assumed to share a
user's logged-in Chrome profile. Authentication, signed-in browser attachment,
permission boundaries, and compatibility need separate trials. Direct account
APIs can offer clearer structured outcomes than visual clicking; browser tools
remain necessary where suitable APIs are unavailable. Account connectors must
be selected and verified individually, including scopes and credential handling.

\textbf{Local Mac tools} require an adapter with explicit permissions and
observable action outcomes. Screenshot/accessibility-based drivers, including
approaches used by OpenClicky, are candidates for investigation. No native driver
has been selected or validated. Broad shell access is not a substitute for a
verified consequential-action boundary.

\subsection{Voice, client, and deployment}
\textbf{LiveKit with Python} is proposed for audio sessions and tool integration
\citep{LiveKitStart,LiveKitTools}. Python aligns with the anticipated agent-side
integration. A TypeScript-only voice backend is an alternative that reduces
language count; the Python choice adds a typed cross-process boundary to test.
A voice tool call does not itself guarantee cancellation of downstream work.

\textbf{TypeScript/React} is proposed for the web client, using semantic controls,
visible text/history, focus management, and full keyboard operation. Familiar web
tooling and iteration speed favor it over an initial Swift-only client. React
does not guarantee accessibility; VoiceOver testing is required. The web interface
and local service support assisted setup followed by independent daily use.

\textbf{ElevenLabs} remains a speech-output candidate through LiveKit's documented
plugin \citep{ElevenLabsPlugin}. Alternatives include other supported speech
providers and native system speech. Selection will compare intelligibility,
interruption behavior, latency, privacy, and actual cost. Recognition and general
reasoning/vision models will likewise use existing providers, with versions and
usage recorded. Local models are an option to investigate, not a commitment to
fully local performance or zero cost.

\textbf{Jev} is optional for structured evaluation or routing
\citep{Jev}. Its documented role does not establish suitability as the primary
conversational or visual-control model. A deterministic policy is the simpler
alternative for mandatory approval rules. Promotional pricing is not a basis
for the project budget.

Deployment remains hybrid: a local Mac, hosted services, or their combination.
Native tasks require an online local component with user-granted permissions.
Hosting the voice/agent services does not remotely confer access to a Mac.
The default topology, secret storage, provider retention, and network requirements
will be documented after a working configuration is demonstrated. A second voice
channel connecting telephone and on-Mac interaction is the leading stretch goal.

\subsection{Engineering and measurement tools}
Git and GitHub provide version history, issues, reviews, and project tracking.
GitHub Actions supplies repeatable checks; hosted CI cannot substitute for Mac
permission tests or participant evaluation. Proposed Python tools are Ruff,
pytest, and pytest-cov; proposed client tools are TypeScript strict checks,
ESLint, Prettier, Vitest, and Testing Library. Final configurations and versions
will be pinned during setup. Coverage is diagnostic, not proof of correctness.

Contract tests will exercise approval, cancellation, reconnection, and task-state
boundaries using controlled workers, followed by real-tool integration trials.
Deterministic tests do not imply deterministic website workflows in the product.
Monotonic timestamps will measure voice response and stop latency; task traces
will measure success, errors, and cost per completed task. A simple instrumented
trace is preferable initially to a complex observability service. The team must
learn the harness protocol, asynchronous audio integration, and accessible focus
management through small tested integration slices.

\section{Use of AI and LLMs}
The product uses existing models to interpret goals and operate tools. The team
also proposes AI assistance for research, code, tests, and documentation, with
human understanding and review. Codex has assisted this draft; other members'
tools and versions remain to be reported. Generated code must pass relevant
checks and review. Sources must be opened and verified, claims distinguished
from assumptions, and participant quotations/results never generated as evidence.
Per-document disclosures record material assistance, affected sections, tool
versions where known, and actual verification \citep{CourseOutline}.

\section{Coding Standard}
Python will follow PEP 8 conventions, type annotations at public boundaries, and
a shared formatter/linter configuration. TypeScript will use strict types and
shared formatting. Small modules separate voice, task control, and tool adapters;
state transitions and error outcomes are explicit. Tests cover behavior and
failure boundaries. Semantic HTML, accessible names, keyboard operation, and
focus behavior are part of review. Credentials and participant data are excluded
from code, fixtures, and ordinary logs. Exceptions to conventions are documented
in the relevant pull request.

\clearpage
\section*{Appendix A: Generative AI Use Disclosure}
\textbf{Tool and use.} Aaron Loh used OpenAI Codex with GPT-6 Astra,
primarily at Medium reasoning and occasionally High (settings reported by
Aaron). Codex helped extract questions from course requirements, challenge
assumptions, research options, and develop ideas through discussion. It also
assisted with drafting and revising text throughout this document, organizing
references, and LaTeX/PDF formatting. Aaron made the judgment calls and reviewed
the work. These are working drafts for further team revision, not evidence of
team-wide agreement or completed participant research.

\textbf{Verification and responsibility.} The assistant checked cited sources,
course requirements, document consistency, and PDF rendering. These checks do
not establish that the proposed system works. Team review, ratification, and
personal reflections remain pending. Before submission, the team will verify
the final content and update this disclosure to reflect subsequent revisions
and any additional AI use. The team remains responsible for the submitted work.

\section*{Appendix B: Reflection}
These answers must come from the members' actual experience, not invented AI
reflections.
\begin{enumerate}
\item Why is it important to create a development plan before starting?
\pending{Members' reflections on this project's planning experience.}
\item What are the advantages and disadvantages of continuous integration and deployment (CI/CD) in your opinion?
\pending{Members' own assessment, including actual experience where available.}
\item What disagreements occurred in this deliverable and how were they resolved?
\pending{An honest account of team discussion; do not claim consensus or disputes
that have not occurred.}
\end{enumerate}

\section*{Appendix C: Team Charter}
\textbf{Status: proposed by Aaron during planning; awaiting ratification by all
four members.} These are proposed working agreements, not evidence of assent.

\subsection*{External goals}
Produce a useful open-source contribution, a credible working prototype informed
by blind users, strong engineering evidence, and a strong course outcome.

\subsection*{Attendance}
\textbf{Expectations.} Attend weekly syncs prepared and on time. Give advance
notice of lateness, early departure, or absence when possible. A written update
and decision input can support flexible attendance; absence alone is not a
failure to contribute.

\textbf{Acceptable excuses.} Illness, emergencies, accessibility needs, and
communicated academic/work conflicts merit rescheduling or reassignment. Repeated
unexplained silence or uncommunicated missed commitments calls for a recovery
conversation, rather than an assumption about personal motivation.

\textbf{Emergency process.} Notify the group when practical, identify work at
risk if able, and let the group redistribute urgent tasks. Personal details need
not be made public.

\subsection*{Accountability and teamwork}
\textbf{Quality.} Weekly commitments specify outputs and acceptance criteria.
Meeting preparation includes reading the agenda and relevant drafts. Contributions
include evidence of checks, understandable documentation, and a reviewable change.
Drafts are ready 48 hours before deadlines. All members must understand and be
able to explain AI-assisted work they submit.

\textbf{Attitude.} Discuss ideas respectfully, make room for disagreement, and
ask before assuming a teammate's intent. Treat participant expertise as evidence;
avoid speaking for blind people without consulting them. Do not disclose personal
participant information to resolve team debates.

\textbf{Staying on track.} Use weekly output commitments and the issue board, not
fixed hours or commit quotas. Two unexplained missed commitments in one month
trigger a private recovery conversation with an explicit revised plan. If the
problem remains unresolved, involve the TA with factual records and an opportunity
for the member to respond. Recognize useful contributions and help remove blockers.

\textbf{Team building.} Include short demos and informal check-ins in the weekly
sync. Recognize progress; any additional social activities are optional and have
not yet been agreed.

\textbf{Decision making.} Seek evidence and run a small experiment where useful,
then vote if discussion does not resolve the question. For a tied technical
choice, the assigned owner decides and records the rationale. Course-policy
questions go to the TA. No individual can invent another member's consent or
participant research results.

\textbf{Ratification.} \pending{All four names, review date, acceptance or amendments,
and initial responsibilities. No signatures or approval dates are implied.}

\clearpage
\bibliographystyle{plainnat}
\bibliography{../../refs/CapstoneSources}
\end{document}
